%!TEX encoding = UTF-8 Unicode
%!TEX root = ../galgas-book.tex

%--------------------------------------------------------------
\chapterLabel{Le type \texttt{array}}{typeArray}
%-------------------------------------------------------------

\tableDesMatieresLocaleDeProfondeurRelative{2}


\section{Déclaration d'un type de tableau}

Un type tableau dont les éléments sont de type \ggsq+@VALUE+ est implicitement déclaré par la construction \ggsq+[@VALUE]+. Le type \ggsq+@VALUE+ est un type quelconque.

\section{Initialisation}

\subsection{Initialiseur sans argument}

L'initialiseur sans argument est implicitement déclaré. Il retourne une liste vide~:

\begin{galgas4}
var [@string] tableau = [@string] ()
\end{galgas4}

Ou, en utilisant l'inférence de type :
\begin{galgas4}
var [@string] tableau = @()
\end{galgas4}


\subsection{Constante litérale}

La construction \ggsq+[@VALUE] [ ... ]+ permet de construire une instance de \ggsq+[@VALUE]+ qui contient les éléments qui apparaissent entre les crochets :

\begin{galgas4}
var [@string] tableau = [@string] ["a", "b", "c"]
\end{galgas4}


Ou, en utilisant l'inférence de type :
\begin{galgas4}
var [@string] tableau = ["a", "b", "c"]
\end{galgas4}


Ou, encode :
\begin{galgas4}
var tableau = [@string] ["a", "b", "c"]
\end{galgas4}

Cas particulier, la liste des éléments peut être vide :
\begin{galgas4}
var [@string] tableau = [@string] []
\end{galgas4}


Ou, en utilisant l'inférence de type :
\begin{galgas4}
var [@string] tableau = []
\end{galgas4}




\section{Opérateurs}

\subsectionLabel{L'instruction \texttt{++=}}{operatorPlusPlusEgalSurListe}

L'instruction \ggsq*++=* ajoute un nouvel élément à la fin du tableau :\\

\begin{galgas4}
var [@string] tableau = []
tableau ++= "a"
tableau ++= "b"
tableau ++= "c"
\end{galgas4}



\subsection{L'instruction \texttt{+=}}

L'instruction \ggsq*cible += expression* ajoute à la fin du tableau \ggsq+cible+ les éléments de la valeur de \ggsq+expression+ :

\begin{galgas4}
var [@string] tableau = []
tableau += ["a"]
let t = [@string] ["b", "c"]
tableau += t
\end{galgas4}





\subsection{L'opérateur \texttt{+}}

L'opérateur \ggsq*+* réalise la concaténation de ses opérandes :

\begin{galgas4}
var [@string] tableau = ["a"] + ["b", "c"]
\end{galgas4}



\section{Fonctions}

\subsection{La fonction \texttt{count}}

\begin{galgas4box}
func [@VALUE].count () -> @u32
\end{galgas4box}

Retourne le nombre d'éléments.


\subsection{La fonction \texttt{first}}

\begin{galgas4box}
func [@VALUE].first () -> @VALUE
\end{galgas4box}

Retourne le premier élément du tableau (à l'indice $0$). Une erreur est déclenché si le tableau est vide



\subsection{La fonction \texttt{last}}

\begin{galgas4box}
func [@VALUE].last () -> @VALUE
\end{galgas4box}

Retourne le dernier élément du tableau. Une erreur est déclenché si le tableau est vide





\subsection{La fonction \texttt{removingFirst}}

\begin{galgas4box}
func [@VALUE].removingFirst (@u32 inIndex) -> [@VALUE]
\end{galgas4box}

Retourne le tableau privé de ses \ggsq*inIndex* éléments. \ggsq*inIndex* doit être inférieur ou égal au nombre d'éléments du tableau. Si égalité, le tableau retourné est vide. Si \ggsq*inIndex* est égal à zéro, l'intégralité du tableau est retourné.






\subsection{La fonction \texttt{sublist}}

\begin{galgas4box}
func [@VALUE].sublist (?from: @u32 inStartIndex, ?count: @u32 inCount) -> [@VALUE]
\end{galgas4box}

Retourne un tableau contenant les \ggsq*inCount* éléments du tableau, à partir de l'indice \ggsq*inStartIndex* éléments.

Si le tableau est vide, un tableau vide est retourné si \ggsq*inStartIndex* et \ggsq*inCount* sont tous les deux égaux à zéro. Si \ggsq*inStartIndex* ou \ggsq*inCount* est supérieur à zéro, une erreur est déclenchée.

Si le tableau n'est pas vide, alors :
\begin{itemize}
  \item \ggsq*inStartIndex* doit être strictement inférieur au nombre d'éléments ;
  \item \ggsq*inStartIndex* + \ggsq*inCount* doit être inférieur ou égal au nombre d'éléments.
\end{itemize}




\section{Méthodes mutantes}




\subsection{La méthode mutante \texttt{append}}

\begin{galgas4box}
mutating proc [@VALUE].append (@VALUE inElement)
\end{galgas4box}

Ajoute un élément à la fin du tableau.


\begin{galgas4}
var [@string] tableau = []
!?tableau.append ("a")
!?tableau.append ("b")
!?tableau.append ("c")
\end{galgas4}



\subsection{La méthode mutante \texttt{removeAll}}

\begin{galgas4box}
mutating proc [@VALUE].removeAll ()
\end{galgas4box}

Supprime tous les éléments du tableau.





\subsection{La méthode mutante \texttt{removeFirst}}

\begin{galgas4box}
mutating proc [@VALUE].removeFirst (!@VALUE outElement)
\end{galgas4box}

Retire le premier élément du tableau, et l'affecte à l'argument de sortie. Une erreur est déclenchée si le tableau est vide.







\subsection{La méthode mutante \texttt{removeLast}}

\begin{galgas4box}
mutating proc [@VALUE].removeLast (!@VALUE outElement)
\end{galgas4box}

Retire le dernier élément du tableau, et l'affecte à l'argument de sortie. Une erreur est déclenchée si le tableau est vide.








\subsection{La méthode mutante \texttt{remove}}

\begin{galgas4box}
mutating proc [@VALUE].remove (!@VALUE outElement, ?at: @u32 inIndex)
\end{galgas4box}

Retire l'élément d'indice \ggsq*inIndex* du tableau, et l'affecte à l'argument de sortie. Une erreur est déclenchée si \ggsq*inIndex* est supérieur ou égal au nombre d'éléments du tableau.









\subsection{La méthode mutante \texttt{insert}}

\begin{galgas4box}
mutating proc [@VALUE].insert (@VALUE inElement, ?at: @u32 inIndex)
\end{galgas4box}

Insère \ggsq*inElement* à l'indice \ggsq*inIndex* du tableau. Une erreur est déclenchée si \ggsq*inIndex* est strictement supérieur au nombre d'éléments du tableau.











\section{Opérateurs d'indexation}



\subsection{Lecture}

\begin{galgas4box}
func [@VALUE] [@u32 inIndex] -> @VALUE
\end{galgas4box}

Retourne l'élement* d'indice \ggsq*inIndex*. Une erreur est déclenchée si \ggsq*inIndex* est supérieur ou égal au nombre d'éléments du tableau.




\subsection{Écriture}

\begin{galgas4box}
mutating func [@VALUE] [@u32 inIndex] = @VALUE
\end{galgas4box}

Affecte la valeur à l'indice \ggsq*inIndex* du tableau. Une erreur est déclenchée si \ggsq*inIndex* est supérieur ou égal au nombre d'éléments du tableau.






















\section{Énumérer un tableau avec l'instruction \texttt{for}}

Par défaut, un tableau est énuméré dans l'ordre des indices croissants :

\begin{galgas4}
let [@string] tableau = ["a", "b", "c"]
for str in tableau do 
  print (str) // Imprime abc
end
println ()
\end{galgas4}

Pour énumérer dans l'ordre inverse, utiliser le délimiteur \ggsq!>! :

\begin{galgas4}
let [@string] tableau = ["a", "b", "c"]
for > str in tableau do 
  print (str) // Imprime cba
end
println ()
\end{galgas4}


